\section{Simulations}
\label{sec:simulation}


In this section, we present the simulation results of TS-MR, Greedy-MR, and other baseline algorithms. 



\subsection{Synthetic datasets}

\begin{figure}[!htp]
	\centering
	\begin{subfigure}[b]{0.32\textwidth}
		\includegraphics[width=\textwidth]{figures/cumu_regret_ex1}
				\caption{example 1}
				\label{fig:ex1}
\end{subfigure}
	\begin{subfigure}[b]{0.32\textwidth}
				\includegraphics[width=\textwidth]{figures/cumu_regret_ex2}
						\caption{example 2}
						\label{fig:ex2}
\end{subfigure}
	\begin{subfigure}[b]{0.32\textwidth}
		\includegraphics[width=\textwidth]{figures/cumu_regret_ex3}
				\caption{example 3}
				\label{fig:ex3}
\end{subfigure}
	\caption{Comparison of the cumulative regret incurred by TS-MR and Greedy-MR, versus the baseline algorithms in Examples 1-3. }
	\label{fig:cumulative_regerts}
\end{figure}

\paragraph{Example 1. Stochastic linear bandit with uniformly and independently distributed actions.} We fix $d = 50$, and sample $\param\sim\text{Unif}(\{\theta \in \mathbb{R}^d|\|\theta\|= 10\}) $ on a sphere with fixed norm. At each time $t$, we generate 100 i.i.d. random actions sampled from $\text{Unif}(\cS_{d-1})$ to form $\cX_t$. This is a basic example of standard stochastic linear bandit problems without any extra structure. We set the threshold  $\mu = 8$ for TS-MR and Greedy-MR. We average simulation results over 100 independent runs. The results are summarized in Figure~\ref{fig:ex1}. TS-Freq shows pessimistic regret due to the inflation of the posterior, while other algorithms in general perform well. 



\paragraph{Example 2. Contextual bandits embedded in the linear bandit problem  \citep{abbasi2013online}.} We fix $d = 50$, and sample $\param\sim\text{Unif}(\{\theta \in \mathbb{R}^d|\|\theta\|= 70\}) $. At each time $t$, we first  generate a random vector $u_t \sim \cN(0, \II_5)$ and let $x_{t,i}\in \RR^{50}$ be the vector whose $i$-th block of size 5 is a  copy of $u_t$, and other components are 0. Then $X_t = \{x_{t,i}\}_{i\in [10] }$ is an action set of size 10, sharing the same feature $u_t$ in different blocks. This problem is equivalent to a 10-armed contextual bandit. We ser $\mu = 12$ for TS-MR and Greedy-MR. Again, we average simulation results over 100 independent runs. The results are summarized in Figure~\ref{fig:ex2}. In this setting, Greedy performs suboptimally due to a lack of exploration for some arms. Nevertheless, Greedy-MR outperforms both Greedy and OFUL by adaptively choosing OFUL actions only when it detects large regret proxy $\hat{\mu}_t$.  



\paragraph{Example 3. Prior mean mismatch \citep{hamidi2020frequentist}.} This is an example in which LinTS is shown to incur linear Bayesian regret. We sample $\theta^\star \sim \cN(m \mathbf{1}_{3d}, \II_{3d})$ and fix the action set $\cX_t = \{0, x_a, x_b\}$ for all $t\in[T]$, where $x_a = - \sum_{i=1}^{d} \frac{e_i}{\sqrt{3d}},  \quad x_b = \sum_{i=11}^{3d} \frac{e_i}{\sqrt{3d}}  - \sum_{i=1}^{d} \frac{e_i}{\sqrt{3d}}$.
It is shown in \cite{hamidi2020frequentist} that, when LinTS takes a wrong prior mean as input,  it has a large probability to choose  $\tilde{x}_2 = 0$, conditioned on $\tilde{x}_1 = x_a$. Note that choosing the zero action brings no information update to LinTS, it suffers a linear Bayesian regret when trying to escape from the zero action. We let $m = 10$ and set $d = 10$, so the problem is a $30$-dimensional linear bandit. We ser $\mu = 12$ for TS-MR and Greedy-MR. We carried out 100 independent runs. The results are summarized in Figure~\ref{fig:ex3}. We see both LinTS and Greedy incur linear regrets as expected, while TS-MR and Greedy-MR, switch to OFUL adaptively to tackle this hard problem  and achieve sublinear regret.



\subsection{Real-world datasets}

We explore the performance of standard POFUL algorithms and the proposed TS-MR and Greedy-MR algorithms on real-world datasets. We consider three datasets for classification tasks on the \href{https://www.openml.org/}{\textsc{OpenML}} platform, \href{https://www.openml.org/d/1466}{Cardiotocography}, \href{https://www.openml.org/d/375}{JapaneseVowels}  and \href{https://www.openml.org/d/36}{Segment}. They focus on healthcare, pattern recognition, and computer vision problems respectively. 

\paragraph{Setup.} We follow the standard approach in the literature that converts classification tasks to contextual bandit problems, and then embeds it into linear bandit problems in the same way as Example 2 in Section 6. Specifically, we regard each class as an action so that at each time, the decision-maker assigns a class to the feature observed and receives a binary reward, namely 1 for assigning the correct class and 0 otherwise, plus a Gaussian noise. 



We plot the cumulative regret (averaged over 100 duplicates) for all algorithms.   Figure~\ref{fig:emp} shows that for all real-world datasets: OFUL and TS-Freq perform poorly due to their conservative exploration; LinTS and Greedy are achieving empirical success even though they don't have theoretical guarantees; TS-MR and Greedy-MR retain the desirable empirical performance of LinTS and Greedy, while enjoying the minimax optimal frequentist regret bound.

 \begin{figure}[htp]
	\centering
	\begin{subfigure}[b]{0.32\textwidth}
		\includegraphics[width=\textwidth]{figures/dataset_cardiotocography.jpg} 
	\caption{Cardiotocography dataset.}
		\label{fig:Cardiotocography}
\end{subfigure}
	\begin{subfigure}[b]{0.32\textwidth}
	\includegraphics[width=\textwidth]{figures/dataset_JapaneseVowels.jpg} 
	\caption{JapaneseVowels dataset.}
		\label{fig:JapaneseVowels}
	\end{subfigure}
 	\begin{subfigure}[b]{0.32\textwidth}
	\includegraphics[width=\textwidth]{figures/dataset_segment.jpg} 
	\caption{Segment dataset.}
		\label{fig:segment}
	\end{subfigure}
	\caption{Cumulative regret of all algorithms on real-world datasets.}
	\label{fig:emp}
\end{figure}








